\documentclass[11pt,letterpaper]{article}
\usepackage[technical-reference]{cornell-notes}
\usepackage{needspace}

\title{Clause 5: Lexical Conventions - Cornell Notes}
\author{}
\date{}
\setDocTitle{Clause 5: Lexical Conventions}
\setDocSubtitle{C++ 2024 Cornell Notes}
\setDocOwner{C++ 2024 Cornell Notes}
\setDocAuthor{}
\setDocDate{}
\setCornellCollection{C++ 2024 Cornell Notes}
\setCornellUnitType{Clause}
\setCornellUnitNumber{5}
\setCornellUnitTitle{Lexical Conventions}
\setCornellCompanionLabel{C++ Technical Reference}
\hypersetup{pdftitle={Clause 5: Lexical Conventions - Cornell Notes},pdfsubject={C++ technical study notes},pdfauthor={}}

\begin{document}
\maketitle
\begin{CornellReferenceOrientationBox}
\textbf{Purpose.} Defines how source characters become preprocessing tokens and tokens through translation phases. It covers character sets, identifiers, keywords, operators, comments, header names, numeric forms, and all literal categories.
\par\textbf{Role.} Provides the lexical substrate consumed by grammar, declarations, expressions, modules, and preprocessing.
\par\textbf{Principal dependencies.} Relies on the translation and conformance model; feeds every syntactic clause and Clause 15.
\par\textbf{Primary audience.} programmers, lexer and preprocessor implementers, and source-analysis tools.
\par\textbf{Major subject groups.} Separate translation; Phases of translation; Character sets; Preprocessing tokens; Alternative tokens; Tokens; Comments; Header names; Preprocessing numbers; Identifiers; Keywords; Operators and punctuators; and related subclauses.
\end{CornellReferenceOrientationBox}
\section{Learning Objectives}
\begin{itemize}[label=\textcolor{CornellReferenceTeal}{$\blacktriangleright$}]
\item Define the governing ideas of Separate translation and state the consequences for a program or implementation.
\item Identify the governing ideas of Phases of translation and state the consequences for a program or implementation.
\item Distinguish the governing ideas of Character sets and state the consequences for a program or implementation.
\item Explain the governing ideas of Preprocessing tokens and state the consequences for a program or implementation.
\item Trace the governing ideas of Alternative tokens and state the consequences for a program or implementation.
\item Apply the governing ideas of Tokens and state the consequences for a program or implementation.
\item Analyze the governing ideas of Comments and state the consequences for a program or implementation.
\item Evaluate the governing ideas of Header names and state the consequences for a program or implementation.
\item Diagnose the governing ideas of Preprocessing numbers and state the consequences for a program or implementation.
\item Compare the governing ideas of Identifiers and state the consequences for a program or implementation.
\end{itemize}
\section{Normative Structure Map}
\small\begin{longtable}{@{}>{\RaggedRight\arraybackslash}p{0.13\textwidth}>{\RaggedRight\arraybackslash}p{0.27\textwidth}>{\RaggedRight\arraybackslash}p{0.19\textwidth}>{\RaggedRight\arraybackslash}p{0.31\textwidth}@{}}
\toprule\textbf{Subclause} & \textbf{Subject} & \textbf{Rule category} & \textbf{Key dependencies / entities}\\\midrule\endfirsthead
\toprule\textbf{Subclause} & \textbf{Subject} & \textbf{Rule category} & \textbf{Key dependencies / entities}\\\midrule\endhead
\bottomrule\endfoot
5.1 & Separate translation & core-language semantics & Relies on the translation and conformance model; feeds every syntactic clause and Clause 15.\\
5.2 & Phases of translation & core-language semantics & Relies on the translation and conformance model; feeds every syntactic clause and Clause 15.\\
5.3 & Character sets & core-language semantics & Relies on the translation and conformance model; feeds every syntactic clause and Clause 15.\\
5.4 & Preprocessing tokens & syntax / lexical rules & Relies on the translation and conformance model; feeds every syntactic clause and Clause 15.\\
5.5 & Alternative tokens & syntax / lexical rules & Relies on the translation and conformance model; feeds every syntactic clause and Clause 15.\\
5.6 & Tokens & syntax / lexical rules & Relies on the translation and conformance model; feeds every syntactic clause and Clause 15.\\
5.7 & Comments & core-language semantics & Relies on the translation and conformance model; feeds every syntactic clause and Clause 15.\\
5.8 & Header names & interface / declarations & Relies on the translation and conformance model; feeds every syntactic clause and Clause 15.\\
5.9 & Preprocessing numbers & core-language semantics & Relies on the translation and conformance model; feeds every syntactic clause and Clause 15.\\
5.10 & Identifiers & core-language semantics & Relies on the translation and conformance model; feeds every syntactic clause and Clause 15.\\
5.11 & Keywords & core-language semantics & Relies on the translation and conformance model; feeds every syntactic clause and Clause 15.\\
5.12 & Operators and punctuators & core-language semantics & Relies on the translation and conformance model; feeds every syntactic clause and Clause 15.\\
5.13 & Literals & syntax / lexical rules & Kinds of literals, Integer literals, Character literals, Floating-point literals, and related subclauses\\
\end{longtable}\normalsize
\begin{CornellReferenceNormativeBox}A heading maps the clause; it does not replace the detailed conditions. For each operation, read requirements, effects, result, exceptions, complexity, remarks, and notes with their stated normative force.\end{CornellReferenceNormativeBox}
\subsection*{Rule Classification Guide}
\small\begin{longtable}{@{}>{\RaggedRight\arraybackslash}p{0.25\textwidth}>{\RaggedRight\arraybackslash}p{0.67\textwidth}@{}}\toprule\textbf{Label or category} & \textbf{How to read it}\\\midrule\endfirsthead\toprule\textbf{Label or category} & \textbf{How to read it (continued)}\\\midrule\endhead\bottomrule\endfoot
Well-formed / ill-formed & Formation is governed by syntax and semantic rules; an ill-formed program does not become well-formed because one implementation accepts it.\\
Diagnosable rule & A violation requires at least one diagnostic unless the wording places the case outside the diagnosable set.\\
No diagnostic required & The program can be ill-formed while the implementation has no diagnostic obligation.\\
Undefined behavior & No execution requirements are imposed; translation success does not establish safety or portability.\\
Unspecified behavior & One of the permitted outcomes may be chosen without documentation.\\
Implementation-defined behavior & The implementation chooses and documents the behavior.\\
Conditionally supported & The implementation may omit support but documents unsupported constructs.\\
\end{longtable}\normalsize
\section{Cornell Cue-and-Notes}
\Needspace{8\baselineskip}
\subsection*{5.1 Separate translation}
\begin{longtable}{@{}>{\RaggedRight\arraybackslash\columncolor{CornellReferenceCueGray}}p{0.28\textwidth}>{\RaggedRight\arraybackslash}p{0.64\textwidth}@{}}
\rowcolor{CornellReferenceNavy}\color{white}\textbf{Cue / recall prompt} & \color{white}\textbf{Notes}\\\endfirsthead
\rowcolor{CornellReferenceNavy}\color{white}\textbf{Cue / recall prompt} & \color{white}\textbf{Notes (continued)}\\\endhead
\arrayrulecolor{CornellReferenceLineGray}
\textbf{What rule applies to Separate translation?}\\[-1mm]{\scriptsize\CornellClauseRef{5.1}} & Defines the core-language rules for Separate translation, including formation, semantic effect, interactions with type and lookup, and any ill-formed or implementation-dependent cases.\\\hline
\end{longtable}
\begin{CornellReferenceSummaryBox}Separate translation supplies the core-language semantics portion of this clause. The key review move is to connect its local wording to the clause-wide model, then classify each condition by normative force before relying on it.\end{CornellReferenceSummaryBox}
\Needspace{8\baselineskip}
\subsection*{5.2 Phases of translation}
\begin{longtable}{@{}>{\RaggedRight\arraybackslash\columncolor{CornellReferenceCueGray}}p{0.28\textwidth}>{\RaggedRight\arraybackslash}p{0.64\textwidth}@{}}
\rowcolor{CornellReferenceNavy}\color{white}\textbf{Cue / recall prompt} & \color{white}\textbf{Notes}\\\endfirsthead
\rowcolor{CornellReferenceNavy}\color{white}\textbf{Cue / recall prompt} & \color{white}\textbf{Notes (continued)}\\\endhead
\arrayrulecolor{CornellReferenceLineGray}
\textbf{What rule applies to Phases of translation?}\\[-1mm]{\scriptsize\CornellClauseRef{5.2}} & Defines the core-language rules for Phases of translation, including formation, semantic effect, interactions with type and lookup, and any ill-formed or implementation-dependent cases.\\\hline
\end{longtable}
\begin{CornellReferenceSummaryBox}Phases of translation supplies the core-language semantics portion of this clause. The key review move is to connect its local wording to the clause-wide model, then classify each condition by normative force before relying on it.\end{CornellReferenceSummaryBox}
\Needspace{8\baselineskip}
\subsection*{5.3 Character sets}
\begin{longtable}{@{}>{\RaggedRight\arraybackslash\columncolor{CornellReferenceCueGray}}p{0.28\textwidth}>{\RaggedRight\arraybackslash}p{0.64\textwidth}@{}}
\rowcolor{CornellReferenceNavy}\color{white}\textbf{Cue / recall prompt} & \color{white}\textbf{Notes}\\\endfirsthead
\rowcolor{CornellReferenceNavy}\color{white}\textbf{Cue / recall prompt} & \color{white}\textbf{Notes (continued)}\\\endhead
\arrayrulecolor{CornellReferenceLineGray}
\textbf{What rule applies to Character sets?}\\[-1mm]{\scriptsize\CornellClauseRef{5.3}} & Specifies text-sequence behavior for Character sets; establish character type, extent, ownership, null termination, encoding assumptions, and invalidation before use.\\\hline
\end{longtable}
\begin{CornellReferenceSummaryBox}Character sets supplies the core-language semantics portion of this clause. The key review move is to connect its local wording to the clause-wide model, then classify each condition by normative force before relying on it.\end{CornellReferenceSummaryBox}
\Needspace{8\baselineskip}
\subsection*{5.4 Preprocessing tokens}
\begin{longtable}{@{}>{\RaggedRight\arraybackslash\columncolor{CornellReferenceCueGray}}p{0.28\textwidth}>{\RaggedRight\arraybackslash}p{0.64\textwidth}@{}}
\rowcolor{CornellReferenceNavy}\color{white}\textbf{Cue / recall prompt} & \color{white}\textbf{Notes}\\\endfirsthead
\rowcolor{CornellReferenceNavy}\color{white}\textbf{Cue / recall prompt} & \color{white}\textbf{Notes (continued)}\\\endhead
\arrayrulecolor{CornellReferenceLineGray}
\textbf{What rule applies to Preprocessing tokens?}\\[-1mm]{\scriptsize\CornellClauseRef{5.4}} & Operates on preprocessing tokens for Preprocessing tokens; expansion, argument substitution, rescanning, and directive context occur before core-language parsing.\\\hline
\end{longtable}
\begin{CornellReferenceSummaryBox}Preprocessing tokens supplies the syntax / lexical rules portion of this clause. The key review move is to connect its local wording to the clause-wide model, then classify each condition by normative force before relying on it.\end{CornellReferenceSummaryBox}
\Needspace{8\baselineskip}
\subsection*{5.5 Alternative tokens}
\begin{longtable}{@{}>{\RaggedRight\arraybackslash\columncolor{CornellReferenceCueGray}}p{0.28\textwidth}>{\RaggedRight\arraybackslash}p{0.64\textwidth}@{}}
\rowcolor{CornellReferenceNavy}\color{white}\textbf{Cue / recall prompt} & \color{white}\textbf{Notes}\\\endfirsthead
\rowcolor{CornellReferenceNavy}\color{white}\textbf{Cue / recall prompt} & \color{white}\textbf{Notes (continued)}\\\endhead
\arrayrulecolor{CornellReferenceLineGray}
\textbf{What rule applies to Alternative tokens?}\\[-1mm]{\scriptsize\CornellClauseRef{5.5}} & Defines the core-language rules for Alternative tokens, including formation, semantic effect, interactions with type and lookup, and any ill-formed or implementation-dependent cases.\\\hline
\end{longtable}
\begin{CornellReferenceSummaryBox}Alternative tokens supplies the syntax / lexical rules portion of this clause. The key review move is to connect its local wording to the clause-wide model, then classify each condition by normative force before relying on it.\end{CornellReferenceSummaryBox}
\Needspace{8\baselineskip}
\subsection*{5.6 Tokens}
\begin{longtable}{@{}>{\RaggedRight\arraybackslash\columncolor{CornellReferenceCueGray}}p{0.28\textwidth}>{\RaggedRight\arraybackslash}p{0.64\textwidth}@{}}
\rowcolor{CornellReferenceNavy}\color{white}\textbf{Cue / recall prompt} & \color{white}\textbf{Notes}\\\endfirsthead
\rowcolor{CornellReferenceNavy}\color{white}\textbf{Cue / recall prompt} & \color{white}\textbf{Notes (continued)}\\\endhead
\arrayrulecolor{CornellReferenceLineGray}
\textbf{What rule applies to Tokens?}\\[-1mm]{\scriptsize\CornellClauseRef{5.6}} & Defines the core-language rules for Tokens, including formation, semantic effect, interactions with type and lookup, and any ill-formed or implementation-dependent cases.\\\hline
\end{longtable}
\begin{CornellReferenceSummaryBox}Tokens supplies the syntax / lexical rules portion of this clause. The key review move is to connect its local wording to the clause-wide model, then classify each condition by normative force before relying on it.\end{CornellReferenceSummaryBox}
\Needspace{8\baselineskip}
\subsection*{5.7 Comments}
\begin{longtable}{@{}>{\RaggedRight\arraybackslash\columncolor{CornellReferenceCueGray}}p{0.28\textwidth}>{\RaggedRight\arraybackslash}p{0.64\textwidth}@{}}
\rowcolor{CornellReferenceNavy}\color{white}\textbf{Cue / recall prompt} & \color{white}\textbf{Notes}\\\endfirsthead
\rowcolor{CornellReferenceNavy}\color{white}\textbf{Cue / recall prompt} & \color{white}\textbf{Notes (continued)}\\\endhead
\arrayrulecolor{CornellReferenceLineGray}
\textbf{What rule applies to Comments?}\\[-1mm]{\scriptsize\CornellClauseRef{5.7}} & Defines the core-language rules for Comments, including formation, semantic effect, interactions with type and lookup, and any ill-formed or implementation-dependent cases.\\\hline
\end{longtable}
\begin{CornellReferenceSummaryBox}Comments supplies the core-language semantics portion of this clause. The key review move is to connect its local wording to the clause-wide model, then classify each condition by normative force before relying on it.\end{CornellReferenceSummaryBox}
\Needspace{8\baselineskip}
\subsection*{5.8 Header names}
\begin{longtable}{@{}>{\RaggedRight\arraybackslash\columncolor{CornellReferenceCueGray}}p{0.28\textwidth}>{\RaggedRight\arraybackslash}p{0.64\textwidth}@{}}
\rowcolor{CornellReferenceNavy}\color{white}\textbf{Cue / recall prompt} & \color{white}\textbf{Notes}\\\endfirsthead
\rowcolor{CornellReferenceNavy}\color{white}\textbf{Cue / recall prompt} & \color{white}\textbf{Notes (continued)}\\\endhead
\arrayrulecolor{CornellReferenceLineGray}
\textbf{What rule applies to Header names?}\\[-1mm]{\scriptsize\CornellClauseRef{5.8}} & Defines the core-language rules for Header names, including formation, semantic effect, interactions with type and lookup, and any ill-formed or implementation-dependent cases.\\\hline
\end{longtable}
\begin{CornellReferenceSummaryBox}Header names supplies the interface / declarations portion of this clause. The key review move is to connect its local wording to the clause-wide model, then classify each condition by normative force before relying on it.\end{CornellReferenceSummaryBox}
\Needspace{8\baselineskip}
\subsection*{5.9 Preprocessing numbers}
\begin{longtable}{@{}>{\RaggedRight\arraybackslash\columncolor{CornellReferenceCueGray}}p{0.28\textwidth}>{\RaggedRight\arraybackslash}p{0.64\textwidth}@{}}
\rowcolor{CornellReferenceNavy}\color{white}\textbf{Cue / recall prompt} & \color{white}\textbf{Notes}\\\endfirsthead
\rowcolor{CornellReferenceNavy}\color{white}\textbf{Cue / recall prompt} & \color{white}\textbf{Notes (continued)}\\\endhead
\arrayrulecolor{CornellReferenceLineGray}
\textbf{What rule applies to Preprocessing numbers?}\\[-1mm]{\scriptsize\CornellClauseRef{5.9}} & Operates on preprocessing tokens for Preprocessing numbers; expansion, argument substitution, rescanning, and directive context occur before core-language parsing.\\\hline
\end{longtable}
\begin{CornellReferenceSummaryBox}Preprocessing numbers supplies the core-language semantics portion of this clause. The key review move is to connect its local wording to the clause-wide model, then classify each condition by normative force before relying on it.\end{CornellReferenceSummaryBox}
\Needspace{8\baselineskip}
\subsection*{5.10 Identifiers}
\begin{longtable}{@{}>{\RaggedRight\arraybackslash\columncolor{CornellReferenceCueGray}}p{0.28\textwidth}>{\RaggedRight\arraybackslash}p{0.64\textwidth}@{}}
\rowcolor{CornellReferenceNavy}\color{white}\textbf{Cue / recall prompt} & \color{white}\textbf{Notes}\\\endfirsthead
\rowcolor{CornellReferenceNavy}\color{white}\textbf{Cue / recall prompt} & \color{white}\textbf{Notes (continued)}\\\endhead
\arrayrulecolor{CornellReferenceLineGray}
\textbf{What rule applies to Identifiers?}\\[-1mm]{\scriptsize\CornellClauseRef{5.10}} & Defines the core-language rules for Identifiers, including formation, semantic effect, interactions with type and lookup, and any ill-formed or implementation-dependent cases.\\\hline
\end{longtable}
\begin{CornellReferenceSummaryBox}Identifiers supplies the core-language semantics portion of this clause. The key review move is to connect its local wording to the clause-wide model, then classify each condition by normative force before relying on it.\end{CornellReferenceSummaryBox}
\Needspace{8\baselineskip}
\subsection*{5.11 Keywords}
\begin{longtable}{@{}>{\RaggedRight\arraybackslash\columncolor{CornellReferenceCueGray}}p{0.28\textwidth}>{\RaggedRight\arraybackslash}p{0.64\textwidth}@{}}
\rowcolor{CornellReferenceNavy}\color{white}\textbf{Cue / recall prompt} & \color{white}\textbf{Notes}\\\endfirsthead
\rowcolor{CornellReferenceNavy}\color{white}\textbf{Cue / recall prompt} & \color{white}\textbf{Notes (continued)}\\\endhead
\arrayrulecolor{CornellReferenceLineGray}
\textbf{What rule applies to Keywords?}\\[-1mm]{\scriptsize\CornellClauseRef{5.11}} & Defines the core-language rules for Keywords, including formation, semantic effect, interactions with type and lookup, and any ill-formed or implementation-dependent cases.\\\hline
\end{longtable}
\begin{CornellReferenceSummaryBox}Keywords supplies the core-language semantics portion of this clause. The key review move is to connect its local wording to the clause-wide model, then classify each condition by normative force before relying on it.\end{CornellReferenceSummaryBox}
\Needspace{8\baselineskip}
\subsection*{5.12 Operators and punctuators}
\begin{longtable}{@{}>{\RaggedRight\arraybackslash\columncolor{CornellReferenceCueGray}}p{0.28\textwidth}>{\RaggedRight\arraybackslash}p{0.64\textwidth}@{}}
\rowcolor{CornellReferenceNavy}\color{white}\textbf{Cue / recall prompt} & \color{white}\textbf{Notes}\\\endfirsthead
\rowcolor{CornellReferenceNavy}\color{white}\textbf{Cue / recall prompt} & \color{white}\textbf{Notes (continued)}\\\endhead
\arrayrulecolor{CornellReferenceLineGray}
\textbf{What rule applies to Operators and punctuators?}\\[-1mm]{\scriptsize\CornellClauseRef{5.12}} & Defines the core-language rules for Operators and punctuators, including formation, semantic effect, interactions with type and lookup, and any ill-formed or implementation-dependent cases.\\\hline
\end{longtable}
\begin{CornellReferenceSummaryBox}Operators and punctuators supplies the core-language semantics portion of this clause. The key review move is to connect its local wording to the clause-wide model, then classify each condition by normative force before relying on it.\end{CornellReferenceSummaryBox}
\Needspace{8\baselineskip}
\subsection*{5.13 Literals}
\begin{longtable}{@{}>{\RaggedRight\arraybackslash\columncolor{CornellReferenceCueGray}}p{0.28\textwidth}>{\RaggedRight\arraybackslash}p{0.64\textwidth}@{}}
\rowcolor{CornellReferenceNavy}\color{white}\textbf{Cue / recall prompt} & \color{white}\textbf{Notes}\\\endfirsthead
\rowcolor{CornellReferenceNavy}\color{white}\textbf{Cue / recall prompt} & \color{white}\textbf{Notes (continued)}\\\endhead
\arrayrulecolor{CornellReferenceLineGray}
\textbf{What rule applies to Kinds of literals?}\\[-1mm]{\scriptsize\CornellClauseRef{5.13.1}} & Defines the core-language rules for Kinds of literals, including formation, semantic effect, interactions with type and lookup, and any ill-formed or implementation-dependent cases.\\\hline
\textbf{What rule applies to Integer literals?}\\[-1mm]{\scriptsize\CornellClauseRef{5.13.2}} & Defines the core-language rules for Integer literals, including formation, semantic effect, interactions with type and lookup, and any ill-formed or implementation-dependent cases.\\\hline
\textbf{What rule applies to Character literals?}\\[-1mm]{\scriptsize\CornellClauseRef{5.13.3}} & Specifies text-sequence behavior for Character literals; establish character type, extent, ownership, null termination, encoding assumptions, and invalidation before use.\\\hline
\textbf{What rule applies to Floating-point literals?}\\[-1mm]{\scriptsize\CornellClauseRef{5.13.4}} & Defines the core-language rules for Floating-point literals, including formation, semantic effect, interactions with type and lookup, and any ill-formed or implementation-dependent cases.\\\hline
\textbf{What rule applies to String literals?}\\[-1mm]{\scriptsize\CornellClauseRef{5.13.5}} & Specifies text-sequence behavior for String literals; establish character type, extent, ownership, null termination, encoding assumptions, and invalidation before use.\\\hline
\textbf{What rule applies to Boolean literals?}\\[-1mm]{\scriptsize\CornellClauseRef{5.13.6}} & Defines the core-language rules for Boolean literals, including formation, semantic effect, interactions with type and lookup, and any ill-formed or implementation-dependent cases.\\\hline
\textbf{What rule applies to Pointer literals?}\\[-1mm]{\scriptsize\CornellClauseRef{5.13.7}} & Defines the core-language rules for Pointer literals, including formation, semantic effect, interactions with type and lookup, and any ill-formed or implementation-dependent cases.\\\hline
\textbf{What rule applies to User-defined literals?}\\[-1mm]{\scriptsize\CornellClauseRef{5.13.8}} & Defines the core-language rules for User-defined literals, including formation, semantic effect, interactions with type and lookup, and any ill-formed or implementation-dependent cases.\\\hline
\end{longtable}
\begin{CornellReferenceSummaryBox}Literals supplies the syntax / lexical rules portion of this clause. The key review move is to connect its local wording to the clause-wide model, then classify each condition by normative force before relying on it.\end{CornellReferenceSummaryBox}
\section{Language-Lawyer and Library-Usage Pitfalls}
\begin{CornellReferencePitfallBox}\begin{itemize}[label=\textcolor{CornellReferenceAmber}{$\blacktriangleright$}]
\item Assuming visual character equality guarantees identifier equivalence.
\item Ignoring the distinction between preprocessing tokens and tokens.
\item Misreading literal type selection after a value exceeds an earlier candidate type.
\item Forgetting encoding and concatenation constraints on string literals.
\item Treating a note, example, or recommended practice as though it imposed a mandatory program rule.
\end{itemize}\end{CornellReferencePitfallBox}
\section{Programmer and Implementation Checklist}
\begin{CornellReferenceChecklistBox}\begin{itemize}[label=$\square$]
\item Tokenize according to translation-phase order.
\item Verify identifier characters, normalization, and reserved forms.
\item Determine each literal's candidate type and encoding.
\item Check raw strings, escape sequences, and suffix rules.
\item Distinguish alternative tokens from unrelated identifiers.
\item For \CornellClauseRef{5.1}, verify separate translation against the precise rule category and all cited dependencies.
\item For \CornellClauseRef{5.2}, verify phases of translation against the precise rule category and all cited dependencies.
\item For \CornellClauseRef{5.3}, verify character sets against the precise rule category and all cited dependencies.
\item For \CornellClauseRef{5.4}, verify preprocessing tokens against the precise rule category and all cited dependencies.
\end{itemize}\end{CornellReferenceChecklistBox}
\section{Clause Synthesis}
\begin{CornellReferenceOrientationBox}Defines how source characters become preprocessing tokens and tokens through translation phases. It covers character sets, identifiers, keywords, operators, comments, header names, numeric forms, and all literal categories. Provides the lexical substrate consumed by grammar, declarations, expressions, modules, and preprocessing. A sound review distinguishes syntax and participation from runtime preconditions, keeps notes and examples non-mandatory, and follows cross-clause dependencies before making a portability claim.\end{CornellReferenceOrientationBox}
\section{Self-Test Questions}
\begin{enumerate}
\item What role does Separate translation play in the clause's overall model?
\item Which normative distinctions must be kept separate when applying Phases of translation?
\item How would you audit a program or implementation that depends on Character sets?
\item Compare Preprocessing tokens with Alternative tokens; what different question does each answer?
\item What role does Alternative tokens play in the clause's overall model?
\item Which normative distinctions must be kept separate when applying Tokens?
\item How would you audit a program or implementation that depends on Comments?
\item Compare Header names with Preprocessing numbers; what different question does each answer?
\item Scenario: a program relies on an unstated implementation behavior while using Preprocessing numbers. What must be checked before calling the program portable?
\item Scenario: a program relies on an unstated implementation behavior while using Identifiers. What must be checked before calling the program portable?
\end{enumerate}
\clearpage\section{Answer Key}
\begin{enumerate}
\item Separate translation is one part of the clause's core-language semantics model. Its role is understood by connecting the local rule in 5.1 to the clause orientation and its dependent concepts.
\item Keep formation and well-formedness separate from runtime preconditions and effects. Also distinguish mandatory wording from notes, implementation choice from undefined behavior, and interface declarations from detailed contracts; see 5.2.
\item Audit Character sets by locating the controlling wording in 5.3, listing inputs and type constraints, proving any preconditions, recording effects and failure behavior, and checking lifetime, invalidation, complexity, and synchronization where applicable.
\item Preprocessing tokens and Alternative tokens occupy different steps or facility groups. Analyze each under its own subclause, then follow the explicit dependency between them rather than transferring one set of guarantees to the other; begin at 5.4.
\item Alternative tokens is one part of the clause's syntax / lexical rules model. Its role is understood by connecting the local rule in 5.5 to the clause orientation and its dependent concepts.
\item Keep formation and well-formedness separate from runtime preconditions and effects. Also distinguish mandatory wording from notes, implementation choice from undefined behavior, and interface declarations from detailed contracts; see 5.6.
\item Audit Comments by locating the controlling wording in 5.7, listing inputs and type constraints, proving any preconditions, recording effects and failure behavior, and checking lifetime, invalidation, complexity, and synchronization where applicable.
\item Header names and Preprocessing numbers occupy different steps or facility groups. Analyze each under its own subclause, then follow the explicit dependency between them rather than transferring one set of guarantees to the other; begin at 5.8.
\item First identify the exact requirement in 5.9, then classify the relied-on behavior as required, unspecified, implementation-defined, conditionally supported, or undefined. Portability requires satisfying the program obligations and avoiding dependence on a choice not guaranteed in the target environments.
\item First identify the exact requirement in 5.10, then classify the relied-on behavior as required, unspecified, implementation-defined, conditionally supported, or undefined. Portability requires satisfying the program obligations and avoiding dependence on a choice not guaranteed in the target environments.
\end{enumerate}
\section{Key-Term Recap}
\begin{longtable}{@{}>{\RaggedRight\arraybackslash}p{0.28\textwidth}>{\RaggedRight\arraybackslash}p{0.64\textwidth}@{}}\toprule\textbf{Term} & \textbf{Working meaning}\\\midrule\endfirsthead\toprule\textbf{Term} & \textbf{Working meaning (continued)}\\\midrule\endhead\bottomrule\endfoot
preprocessing token & A lexical unit recognized before full token classification.\\
token & A lexical element consumed by the syntactic grammar.\\
identifier & A name-forming token subject to character and normalization rules.\\
literal & A source notation denoting a value or string-like object.\\
user-defined literal & A literal form whose suffix participates in user-defined interpretation.\\
translation phase & An ordered conceptual stage that transforms a source file into a translation unit.\\
\end{longtable}
\CornellTakeaway{Lexical classification fixes the tokens and literal meanings on which every later syntactic and semantic rule depends.}
\end{document}
